% ============================================================
%  N.1.3  Partial and Full Fusion
% ============================================================
\subsection{Partial and Full Fusion}
\label{sec:partial-full-fusion}

As discussed in Section~\ref{sec:workloads}, full fusion of an entire
workload can be very demanding even for a fusion-specialised
accelerator, because the on-chip memory must simultaneously
accommodate the weights, activations, and intermediate data of all
fused layers.
To capture fusion benefits without requiring an architecture that can
host the full workload at once, we also evaluate \emph{partial fusion},
in which only a small group of consecutive layers (typically two or
three) is fused into a single segment.
Each segment is executed independently, and the total energy and latency
are obtained by summing over all segments:
%
\begin{equation}
  E_{\text{total}} = \sum_{i} E_{i},
  \qquad
  L_{\text{total}} = \sum_{i} L_{i},
  \qquad
  \text{EDP} = E_{\text{total}} \times L_{\text{total}}.
  \label{eq:sum-partial}
\end{equation}
%
Note that the EDP is the product of the summed energy and summed
latency, not the sum of individual per-segment EDPs, since the
segments are executed sequentially.

\paragraph{Example: FSRCNN.}
FSRCNN has 8~convolutional layers.
Under 2-layer partial fusion the workload is split into four segments,
each fusing two consecutive layers: L0\,+\,L1, L2\,+\,L3, L4\,+\,L5,
and L6\,+\,L7.
Each segment is mapped and executed independently on the same
architecture, and the results are aggregated with
Equation~\eqref{eq:sum-partial}.
In contrast, under full fusion all 8~layers are fused into a single
execution.

\paragraph{Fusion granularities.}
Table~\ref{tab:fusion-levels} lists the partial-fusion granularities
available for each workload.
For FSRCNN and MC-CNN the intermediate level is 2-layer fusion;
for VGG16 it follows the network's block boundaries (pairs of layers
within the same spatial resolution), and for ResNet18 both 2-layer and
stage-level (residual block) groupings are defined.

\begin{table}[htbp]
  \centering
  \caption{Fusion granularities per workload.}
  \label{tab:fusion-levels}
  \small
  \begin{tabular}{l c c l}
    \toprule
    \textbf{Workload}
      & \textbf{Fused layers}
      & \textbf{Segments}
      & \textbf{Partial-fusion level} \\
    \midrule
    FSRCNN   &  8 & 4 & 2-layer \\
    MC-CNN   &  4 & 2 & 2-layer \\
    VGG16    & 13 & 8 & Block (2-layer blocks + solos) \\
    ResNet18 & 17 & 8 & 2-layer \\
    \bottomrule
  \end{tabular}
\end{table}

% ------------------------------------------------------------------
%  Scenario A and Scenario B
% ------------------------------------------------------------------
\subsubsection{Evaluation Scenarios}
\label{sec:scenarios}

Because the architecture must be sized differently depending on the
fusion level it supports, we define two evaluation scenarios:

\begin{description}
  \item[Scenario~A (full-sized architecture).]
    The on-chip memories and registers are dimensioned to support
    \emph{full fusion} of the entire workload, i.e.\ they are large
    enough to hold the weights and intermediate data of all fused layers
    simultaneously.
    Three configurations are compared on this same full-sized
    architecture:
    (i)~\emph{full fusion}, where all layers are fused into one
    execution;
    (ii)~$\Sigma$\,\emph{partial fusion}, where partial-fusion
    segments are executed independently and their results summed; and
    (iii)~$\Sigma$\,\emph{singles}, where every layer is executed
    individually.
    This provides a fair comparison, as all three strategies use
    identical hardware.

  \item[Scenario~B (partial-sized architecture).]
    The on-chip memories and registers are dimensioned only for the
    largest partial-fusion segment, resulting in a \emph{smaller and
    less energy-intensive} architecture than Scenario~A.
    Full fusion is not feasible on this reduced hardware, so only
    $\Sigma$\,partial fusion and $\Sigma$\,singles are compared.
    This scenario answers a practical question: if the silicon budget
    does not permit the large memories required for full fusion, is
    partial fusion still beneficial compared with single-layer
    execution on the same smaller chip?
\end{description}

\noindent
Table~\ref{tab:scenarios-example} illustrates the memory sizing
difference between the two scenarios for selected configurations.

\begin{table}[htbp]
  \centering
  \caption{Example architecture sizing for Scenario~A and Scenario~B.}
  \label{tab:scenarios-example}
  \small
  \begin{tabular}{l l l l}
    \toprule
    \textbf{Configuration}
      & \textbf{Scenario~A}
      & \textbf{Scenario~B}
      & \textbf{Difference} \\
    \midrule
    Eyeriss, ResNet18 (512$\times$32)
      & WReg\,=\,902
      & WReg\,=\,384
      & $2.3\times$ smaller \\
    Eyeriss, VGG16 (512$\times$32)
      & WReg\,=\,1\,200
      & WReg\,=\,576
      & $2.1\times$ smaller \\
    DepFiN, FSRCNN
      & FMEM\,=\,576\,KB
      & FMEM\,=\,248\,KB
      & $2.3\times$ smaller \\
    DepFiN, ResNet18
      & WMEM\,=\,10\,738\,KB
      & WMEM\,=\,4\,608\,KB
      & $2.3\times$ smaller \\
    \bottomrule
  \end{tabular}
\end{table}
